\section{Design and Methodology}

\subsection{Design principles}
Two rules shape the design. (1) \emph{The LLM never sees the raw dump}, only structured findings. This
shrinks the attack surface and keeps every statement traceable to a finding id. (2) \emph{Facts come from
code, not from the LLM.} Cipher confidence, key strength and the risk score are computed
deterministically; the LLM only writes them up and a validator checks that it did so honestly.

\subsection{Block diagram}
Fig.~\ref{fig:block} shows the six-stage pipeline. Stages 1 to 5 write findings to SQLite; stage 6 reads
them back and produces either a rule-based report (always available) or a defended LLM report.

\begin{figure*}[t]
\centering
\resizebox{\textwidth}{!}{%
\begin{tikzpicture}[stage/.style={blk, text width=62mm, minimum height=9mm}]
  % stages
  \node[stage] (s1) at (6.0,0)     {\textbf{1 Memory}: Volatility3 processes, command lines, connections, malfind, USB history};
  \node[stage] (s2) at (6.0,-1.3)  {\textbf{2 Static}: strings, PE headers, imports, section entropy, YARA};
  \node[stage, fill=green!12] (s3) at (6.0,-2.6) {\textbf{3 Crypto (core)}: entropy, cipher ID with confidence, RSA extraction + weak keys};
  \node[stage] (s4) at (6.0,-3.9)  {\textbf{4 Credentials}: JWT, OAuth/API tokens, OTP secrets (redacted)};
  \node[stage] (s5) at (6.0,-5.2)  {\textbf{5 Network}: pyshark, TLS version and suite, DNS, SNI};
  % inputs
  \node[io, text width=17mm] (bin)  at (0,-1.3) {Binary\\(optional)};
  \node[io, text width=20mm] (dump) at (0,-2.6) {Memory dump\\(.raw/.dmp)};
  \node[io, text width=17mm] (pcap) at (0,-5.2) {pcap\\(optional)};
  % dump bus with a hop where the binary line crosses it
  \draw (dump.east) -- (1.8,-2.6);
  \draw (1.8,0) -- (1.8,-3.9);
  \draw[->] (1.8,0) -- (s1.west);
  \draw[->] (1.8,-2.6) -- (s3.west);
  \draw[->] (1.8,-3.9) -- (s4.west);
  \draw[->] (bin.east) -- (1.65,-1.3) arc(180:0:0.15) -- (s2.west);
  \draw[->] (pcap.east) -- (s5.west);
  \draw[->] (s1.south) -- node[right, font=\scriptsize]{extracted RWX regions} (s2.north);
  % collector bus into SQLite
  \foreach \s in {s1,s2,s3,s4,s5} \draw (\s.east) -- (\s.east -| 9.6,0);
  \draw (9.6,0) -- (9.6,-5.2);
  \node[db] (db) at (10.9,-2.6) {SQLite\\findings\\F1, F2, ...};
  \draw[->] (9.6,-2.6) -- (db.west);
  % reporting
  \node[blk, text width=17mm] (risk) at (13.1,-1.6) {Rule-based\\risk 0 to 100};
  \node[blk, text width=17mm] (rule) at (13.1,-3.6) {6a Rule-based\\report};
  \node[llm, text width=30mm] (llm)  at (16.6,-1.6) {\textbf{6b LLM report}\\sanitise, spotlight, LLM,\\validate, retry once};
  \node[io, text width=22mm] (out)   at (16.6,-3.6) {results/ +\\dashboard};
  \draw[->] (db) -- (risk);
  \draw[->] (db) -- (rule);
  \draw[->] (risk) -- (llm);
  \draw[->] (db.north) -- (10.9,-0.35) -- node[above, font=\scriptsize]{findings only, never the raw dump} (16.6,-0.35) -- (llm.north);
  \draw[->] (llm) -- (out);
  \draw[->] (rule) -- (out);
\end{tikzpicture}}
\caption{TriageGuard block diagram. Every stage records numbered findings; the LLM sees findings only.}
\label{fig:block}
\end{figure*}

\subsection{Architecture}
The code is a Python package (\texttt{triageguard/}) with one module per stage, a SQLite store, a report
layer and a command-line front end (Fig.~\ref{fig:arch}). \texttt{pipeline.py} orchestrates the stages;
each stage is independent, so a missing input (no Windows kernel, no pcap, no API key) skips that stage
instead of failing the run. A local dashboard (standard library HTTP server, localhost only, offline)
reads the exported \texttt{results/} folder.

\begin{figure}[t]
\centering
\resizebox{\columnwidth}{!}{%
\begin{tikzpicture}[every node/.style={font=\scriptsize}]
  \node[blk, text width=80mm, fill=gray!10] (ui) at (0,0) {\textbf{Interface}: \texttt{cli.py} (analyze, crypto, report, redteam, export, dashboard), \texttt{dashboard.py}};
  \node[blk, text width=80mm, fill=yellow!10] (orch) at (0,-1.0) {\textbf{Orchestration}: \texttt{pipeline.py} (stage order, severities, skipping)};
  \node[blk, text width=13mm] (m)  at (-3.3,-2.25) {\texttt{memory}\\\texttt{static}};
  \node[blk, text width=28mm, fill=green!12] (c) at (-0.65,-2.25) {\texttt{crypto/}: entropy, identify,\\signatures, rsa\_weak};
  \node[blk, text width=16mm] (cr) at (2.05,-2.25) {\texttt{credentials}\\\texttt{network}};
  \node[blk, text width=11mm, fill=red!8] (rt) at (3.75,-2.25) {\texttt{redteam}};
  \node[draw, dashed, fit=(m)(c)(cr)(rt), inner sep=1mm] (an) {};
  \node[blk, text width=80mm, fill=red!8] (rep) at (0,-3.5) {\textbf{Reporting}: \texttt{report.py} (risk, rule report, LLM client), \texttt{defences.py} (D1 to D4)};
  \node[db] (st) at (-2.0,-4.8) {\texttt{store.py}\\SQLite};
  \node[ext] (ext) at (2.0,-4.8) {External: Volatility3, pefile,\\YARA, tshark, Groq API};
  \draw[->] (ui) -- (orch);
  \draw[->] (orch) -- (an);
  \draw[->] (an) -- (rep);
  \draw[->] (rep.south -| st) -- (st);
  \draw[->, dashed] (rep.south -| ext) -- (ext);
\end{tikzpicture}}
\caption{Layered architecture of the Python package (analysis modules in the dashed box).}
\label{fig:arch}
\end{figure}

\subsection{Data flow (DFD level 0)}
Fig.~\ref{fig:dfd0} shows the context diagram. The level-1 DFD and UML diagrams are in the Appendix.

\begin{figure}[t]
\centering
\resizebox{\columnwidth}{!}{%
\begin{tikzpicture}
  \node[proc, minimum size=19mm] (p) at (0,0) {0\\TriageGuard};
  \node[ext, text width=17mm] (an)   at (-3.6,1.7)  {Analyst};
  \node[ext, text width=17mm] (ds)   at (-3.6,-1.7) {Public dump\\+ pcap};
  \node[ext, text width=19mm] (vol)  at (3.6,1.7)   {Volatility3 +\\MS symbols};
  \node[ext, text width=19mm] (groq) at (3.6,-1.7)  {LLM API\\(Groq)};
  \draw[flow] (an) to[bend left=12] node[above, sloped]{command} (p);
  \draw[flow] (p) to[bend left=12] node[below, sloped]{report, dashboard} (an);
  \draw[flow] (ds) -- node[below, sloped]{image, capture} (p);
  \draw[flow] (p) to[bend left=12] node[above, sloped, pos=0.6]{dump path} (vol);
  \draw[flow] (vol) to[bend left=12] node[below, sloped, pos=0.35]{JSON rows} (p);
  \draw[flow] (p) to[bend left=12] node[above, sloped, pos=0.6]{sanitised findings} (groq);
  \draw[flow] (groq) to[bend left=12] node[below, sloped, pos=0.4]{report text} (p);
\end{tikzpicture}}
\caption{DFD level 0 (context diagram).}
\label{fig:dfd0}
\end{figure}

\subsection{Database design (ER diagram)}
All evidence lives in two tables (Fig.~\ref{fig:er}). A \texttt{run} is one analysis of one dump; it owns
many \texttt{findings}. Each finding has a citation id (\texttt{ref}, e.g.\ F12) that the LLM must cite and the
validator and dashboard resolve. Stage-specific detail is a JSON column, so new checks need no schema change.

\begin{figure}[t]
\centering
\begin{tikzpicture}[
  ent/.style={draw, rectangle split, rectangle split parts=2, align=left, font=\scriptsize,
              rectangle split part fill={blue!15,white}, inner sep=2pt, text width=33mm}]
  \node[ent] (runs) {\textbf{runs}\nodepart{second}
    \underline{id} INTEGER PK\\ started TEXT (UTC ISO)\\ dump TEXT\\ pcap TEXT};
  \node[ent, right=9mm of runs] (f) {\textbf{findings}\nodepart{second}
    \underline{id} INTEGER PK\\ run\_id INTEGER FK\\ ref TEXT (F1, F2, ...)\\ module TEXT\\ kind TEXT\\
    title TEXT\\ severity TEXT\\ source TEXT\\ location TEXT\\ data TEXT (JSON)};
  \draw (runs.east |- runs.north) ++(0,-6mm) coordinate (a) -- node[pos=0.12,above,font=\scriptsize]{1} node[pos=0.88,above,font=\scriptsize]{N} node[midway,below,font=\scriptsize]{has} (a -| f.west);
  \node[font=\scriptsize, below=1mm of runs, text width=33mm, align=left]
    {module $\in$ \{memory, static, crypto, rsa, credentials, network\}\\[1mm]
     severity $\in$ \{info, low, medium, high, critical\}\\[1mm]
     one run \textit{has} many findings};
\end{tikzpicture}
\caption{ER diagram of the SQLite schema (\texttt{store.py}).}
\label{fig:er}
\end{figure}

\subsection{Modules, algorithms and justification}

\textit{Memory (Volatility3).} The standard open-source framework; it auto-detects the Windows build and
fetches symbols. Raw \texttt{malfind} hits are graded: an \texttt{MZ} header in a region is critical, known
JIT processes and zero-filled regions are low, identical code in three or more processes is a system-wide
component (medium), and unique code in one process is high. Suspicious command lines (encrypt, ransom,
vssadmin, bcdedit) and shells with unusual parents are flagged.

\textit{Static (pefile, YARA).} Applied to extracted regions and supplied binaries: SHA-256, entropy,
ASCII and UTF-16 strings, PE imports and section entropy, and YARA rules.

\textit{Entropy.} Shannon entropy $H=-\sum_b p_b\log_2 p_b$ over a 4\,KB sliding window; windows with
$H\ge7.2$ bits/byte are merged into encrypted or packed regions.

\textit{Cipher identification.} 56 signatures (S-boxes, T-tables, round constants, DES tables, the ChaCha
``expand 32-byte k'' constant, GHASH tables, curve parameters, key-blob headers and API names, in ASCII
and UTF-16), each with a weight $w_i\in(0,1)$. For algorithm $A$ with matched signatures $S_A$:
\begin{equation}
  \mathrm{conf}(A) = 1 - \prod_{i\in S_A}(1-w_i).
\end{equation}
One strong hit counts a lot, weak hits accumulate, and confidence never reaches 100\%. Each result also
states its \emph{basis}: implementation constants (the cipher is implemented here) or API/string references
only (merely named), because almost every Windows process names AES.

\textit{RSA weak keys.} Keys are carved from CAPI blobs, CNG blobs, DER SubjectPublicKeyInfo and PEM, then
tested for modulus size ($\le$512 critical, $<$1024 high, $<$2048 medium), small exponent, shared modulus,
shared prime by pairwise $\gcd(n_i,n_j)$, close primes by Fermat factorisation, and small factors (trial
division, Pollard rho). The strength score is a base of 100 ($\ge$3072 bits), 90 ($\ge$2048) or 60, minus a
penalty for the worst issue (critical 100, high 50, medium 25, low 10), clamped to $[0,100]$. A modulus with a
tiny factor is treated as corrupt memory, not a key.

\textit{Credentials.} JWTs are decoded and checked (\texttt{alg=none} high, HS256 noted, expiry);
OTP secrets under 128 bits (the RFC~4226 minimum) are flagged; all secrets are stored redacted with a short hash.

\textit{Network (pyshark).} Uses Wireshark's dissectors. Only the server's chosen suite and version are
read; RC4, DES/3DES, NULL/export, static RSA key exchange, CBC, SSL~3.0 and TLS~1.0/1.1 are flagged.

\textit{Risk score.} Points per severity $p=\{25,12,5,2,0\}$ for critical to info, with each count capped
$c=\{4,4,4,5\}$ so low-severity noise cannot add up to critical:
\begin{equation}
  \mathrm{risk} = \min\Big(100,\ \sum_{s} p_s\cdot\min(n_s, c_s)\Big).
\end{equation}

\textit{LLM report and defences.} Groq's free tier (\texttt{openai/gpt-oss-120b}) through an
OpenAI-compatible API, so any endpoint (e.g.\ local Ollama) can be swapped in. Four defences:
(D1) \emph{sanitise} instruction-like text and forged citations, each hit becoming a high
\texttt{prompt\_injection} finding; (D2) \emph{spotlighting}~\cite{hines2024spotlighting}: findings are marked as
untrusted data between delimiters; (D3) \emph{citation check}: $\ge$90\% of claim lines cited and every id
must exist; (D4) \emph{rule cross-check}: stated risk within 20 points of ours, every critical/high finding
mentioned, every named algorithm present, dismissive or dangerous advice flagged. A failing report gets
one retry with the list of problems; if it still fails it ships with a ``Validation notes'' section.
The sequence is shown in Fig.~\ref{fig:seq} (Appendix).

\textit{Why SQLite and Python.} A single-file database needs no server, is easy to export, and gives each
finding a stable id. Python has mature bindings for every tool used (Volatility3 is itself Python).

\subsection{Work plan}
Table~\ref{tab:plan} lists the plan with its status at this evaluation.

\begin{table}[t]
\centering
\caption{Work plan and status}
\label{tab:plan}
\footnotesize
\begin{tabularx}{\columnwidth}{@{}lXc@{}}
\toprule
Phase & Work item & Status\\
\midrule
P1 & Pipeline skeleton, SQLite store, CLI & \done\\
P2 & Crypto module: entropy, cipher ID (widened scope), RSA weak keys & \done\\
P3 & Credentials, network, rule report, Groq LLM report & \done\\
P4 & Real dumps (Win11, Win10 + pcap, Win7), severity tuning & \done\\
P5 & Red-team (6 attacks $\times$ 2 fields) and four defences, before/after run & \done\\
P6 & Dashboard and automated tests (28) & \done\\
\midrule
P7 & Ground-truth test in our own VM (AES, ChaCha20, RSA, EC keys we generate) & \todo\\
P8 & More public dumps (ransomware-style Win10/11, Windows Server) & \todo\\
P9 & Stronger red-team: multi-field and trigger-free attacks, more trials, second LLM & \todo\\
P10 & Final report, literature survey (Overleaf), PQC discussion & \todo\\
\bottomrule
\end{tabularx}
\end{table}
